\section{Introduction}
\label{sec:intro}

In class-incremental learning (CIL), a model learns a sequence of tasks with
disjoint classes and must classify all classes seen so far. In the
exemplar-free setting, no image of a past task may be kept. RSIAT~\cite{rsiat}
is a strong recent method for this setting. It trains one shared adapter on a
frozen ViT-B/16. After every task it stores the mean and covariance of each
new class, and it uses them to re-train the classifier on sampled
pseudo-features (\emph{classifier alignment}, CA~\cite{zhang2023slca}). The
statistics are therefore already in memory, but at inference time only the
re-trained classifier is used.

\paragraph{Goal.}
This work had two aims. The first was to improve RSIAT with quantum
machine-learning components: a projected quantum kernel~\cite{huang2021power}
as a learned metric, and the pretty-good measurement
(PGM)~\cite{hausladen1994pgm} as a classifier over class states. The second
was to decide, with measurements rather than single runs, whether any change
truly helps. Training randomness alone moves RSIAT by $0.14$ points of
average incremental accuracy on ImageNet-R and by $0.38$ on ImageNet-A, which
is as large as many reported improvements.

\paragraph{Idea.}
A diagnosis of RSIAT (Section~\ref{sec:gap}) shows that the remaining errors
are mostly confusions between classes of different tasks, concentrated on
the oldest classes. It also shows that better statistics alone do not fix
them. However, the stored statistics define a \emph{generative} classifier
that errs on different images than the CA classifier. We therefore add
the two scores at inference time, after putting them on a common scale. The
fusion needs no training, no new parameters and no exemplars. The generative
score can be computed classically (ridge regression on the stored moments)
or as a quantum measurement (a PGM on class density matrices), which gives a
direct test of the quantum component.

\paragraph{Contributions.}
\begin{itemize}
  \item \emph{Statistics--classifier fusion} (SCF), a training-free inference
    rule for exemplar-free CIL methods that store class statistics
    (Section~\ref{sec:method}). On four benchmarks with three seeds, it
    raises final accuracy on all four. It raises average accuracy on three
    of them, and not on CUB-200.
  \item A controlled evaluation of five quantum placements, each against a
    matched classical counterpart. None helps, and we explain why
    (Section~\ref{sec:quantum}).
  \item A diagnosis of where RSIAT loses accuracy: cross-task confusion of
    the oldest classes, which oracle statistics only move to newer classes
    (Section~\ref{sec:gap}).
  \item A measurement protocol for CIL: paired seeds, all inference heads
    evaluated on the same network with bit-identity checks, and
    pre-registered decision rules, including the tests that failed
    (Section~\ref{sec:protocol}).
\end{itemize}
